%\vspace{-2mm}
\section{Conclusion and Future Work} \label{sec:future}
%\vspace{-3mm}
% In this resource paper, we create a new corpus \ilpcsr\ that brings together PCR and LSI tasks for the first time. We experimented with a wide variety of methods for each of the task. Our experiments show that ensemble of lexical (BM-25) and deep semantic method (GNN) perform the best. However, modeling both the tasks in a multi-task setting degrades the performance a bit for one of the task. In future, we plan to work more with the ensembling techniques, including devising fine-tuning approaches to incorporate BM25 scores. We also plan to use multiple representations/features to construct the models, since we believe this will help in tuning statutes and precedents together. 
We introduce \ilpcsr{}, the first corpus that enables parallel retrieval of legal statutes and precedents. Our experiments yield three central insights. First, \textit{ensembles} of lexical and semantic methods (Para-GNN (summaries) + BM25 with dynamic $\alpha$) are the strongest first-stage retrievers. Second, multi-task training degrades performance, while \textit{transfer learning} between LSR and PCR reliably improves results. Third, \textit{LLM-based re-ranking}, with our two-stage cross-conditioning prompt, produces the best overall performance and demonstrates that LSR and PCR can enhance each other. In future, we plan to explore re-ranking with diverse LLM families and prompt designs. 
We also plan to explore finer-grained supervision (e.g., para or rhetorical-role level) to better capture legal relevance. 

